% Teil II: der durchgerechnete Weg mit Begründung, Kontrolle und Merksatz.

\subsection*{Gegeben und gesucht}

Gegeben sind die Masse $m = \qty{0.20}{\kilo\gram}$, die Federkonstante
$k = \qty{5.0}{\newton\per\metre}$ und die Amplitude $A = \qty{4.0}{\centi\metre}$.
Gesucht sind Periodendauer $T$, Frequenz $f$, die größte Geschwindigkeit
$v_{\max}$ und die größte Beschleunigung $a_{\max}$.

\subsection*{Ansatz}

Die Feder liefert die einzige rücktreibende Kraft, Reibung gibt es nicht, also
ist das ein harmonischer Oszillator. Der Weg über das zweite Newtonsche Gesetz
liefert $\omega$ direkt aus $k$ und $m$; alle weiteren Größen hängen nur noch
von $\omega$ und $A$ ab.

\subsection*{Rechnung}

Kraftansatz und Umformen auf die Oszillatorform:
\begin{align}
  m\,\ddot{x} &= -k\,x
  &\Longrightarrow\qquad
  \ddot{x} + \frac{k}{m}\,x &= 0 ,
\end{align}
also mit $\omega^2 = k/m$:
\begin{equation}
  \omega = \sqrt{\frac{k}{m}}
         = \sqrt{\frac{\qty{5.0}{\newton\per\metre}}{\qty{0.20}{\kilo\gram}}}
         = \sqrt{\qty{25}{\per\second\squared}}
         = \qty{5.0}{\per\second} .
\end{equation}

Daraus Periodendauer und Frequenz:
\begin{equation}
  T = \frac{2\pi}{\omega}
    = \frac{2\pi}{\qty{5.0}{\per\second}}
    = \qty{1.3}{\second},
  \qquad
  f = \frac{1}{T} = \qty{0.80}{\hertz} .
\end{equation}

Die Geschwindigkeit ist die erste Ableitung von $x(t) = A\cos(\omega t+\varphi)$,
ihr Betrag wird maximal, wenn der Sinus $\pm 1$ ist:
\begin{equation}
  v_{\max} = \omega A = \qty{5.0}{\per\second}\cdot\qty{0.040}{\metre}
           = \qty{0.20}{\metre\per\second} .
\end{equation}

Für die größte Beschleunigung genügt es, den Scheitelpunkt zu betrachten, dort
ist die Federkraft am größten:
\begin{equation}
  a_{\max} = \omega^2 A
           = \left(\qty{5.0}{\per\second}\right)^2\cdot\qty{0.040}{\metre}
           = \qty{1.0}{\metre\per\second\squared} .
\end{equation}

\begin{ergebnis}
\begin{equation*}
  \omega = \qty{5.0}{\per\second},\quad
  T = \qty{1.3}{\second},\quad
  f = \qty{0.80}{\hertz},\quad
  v_{\max} = \qty{0.20}{\metre\per\second},\quad
  a_{\max} = \qty{1.0}{\metre\per\second\squared}.
\end{equation*}
\end{ergebnis}

\begin{kontrolle}
\textbf{Einheiten:} $\sqrt{\si{\newton\per\metre}/\si{\kilo\gram}}
= \sqrt{\si{\per\second\squared}} = \si{\per\second}$, das passt zu $\omega$.
\textbf{Querprobe über die Energie:} die Gesamtenergie muss in beiden
Schreibweisen dieselbe sein,
\[
  \tfrac{1}{2}kA^2 = \tfrac{1}{2}\cdot\qty{5.0}{\newton\per\metre}\cdot(\qty{0.040}{\metre})^2
  = \qty{4.0}{\milli\joule},
  \qquad
  \tfrac{1}{2}mv_{\max}^2 = \tfrac{1}{2}\cdot\qty{0.20}{\kilo\gram}\cdot(\qty{0.20}{\metre\per\second})^2
  = \qty{4.0}{\milli\joule}.
\]
\textbf{Grenzfall:} für $k\to 0$ geht $T\to\infty$, eine weiche Feder schwingt
langsam; für große $k$ wird $T$ klein. Beides ist plausibel.
\textbf{Größenordnung:} eine Schwingung dauert gut eine Sekunde, die
Maximalgeschwindigkeit ist kleiner als ein halber Meter pro Sekunde.
\end{kontrolle}

\begin{merksatz}
Bei harmonischen Schwingungen kommt es fast nie auf die Zeitfunktion an, sondern
auf $\omega = \sqrt{k/m}$, die Amplitude und die beiden Extremfälle
$v_{\max}=\omega A$ und $a_{\max}=\omega^2 A$. Wer das einmal aufgeschrieben hat,
kann jede Frage zu Frequenz, Periode, Geschwindigkeit und Beschleunigung ohne
Umweg über die Differentialgleichung beantworten.
\end{merksatz}
